\subsection{Caso 10. Mensajería XMPP: interceptación y errores de TLS}
\label{caso:xmpp-tls}

\begin{situacion}
Se investiga si la comunicación entre clientes y servidor de mensajería está correctamente protegida. XMPP se especifica en el RFC~6120 y TLS~1.3 en el RFC~8446 \parencite{rfc6120,rfc8446}.
\end{situacion}

\paragraph{Procedimiento.}
\begin{itemize}[leftmargin=*, topsep=2pt, itemsep=2pt]
	\item En un servidor XMPP de laboratorio (por ejemplo, Openfire o ejabberd), revisar la configuración TLS, los certificados, los nombres, la vigencia y la cadena de confianza.
	\item Capturar metadatos de conexión en una prueba autorizada y verificar que el contenido de una sesión cifrada no se trate como texto legible.
	\item Revisar los registros del servidor en busca de errores de negociación, certificados inválidos o conexiones inesperadas.
\end{itemize}

\paragraph{Indicadores que deben buscarse.}
Certificados vencidos o no confiables, negociación insegura, servicios de administración expuestos o patrones de conexión que no coinciden con el uso esperado.

\paragraph{Mitigación y defensa.}
Exigir TLS robusto y validación de certificados; limitar las interfaces administrativas; mantener el servidor actualizado; aplicar cifrado de extremo a extremo (E2EE) cuando el modelo de amenaza lo requiera.

\paragraph{Evidencias que se entregan.}
Checklist TLS, eventos anonimizados, evidencia de configuración y plan de corrección.

\paragraph{Preguntas de análisis.}
\begin{enumerate}[leftmargin=*, topsep=2pt, itemsep=2pt]
	\item ¿Qué protege TLS y qué no protege?
	\item ¿En qué se diferencia el cifrado de transporte del cifrado de extremo a extremo?
\end{enumerate}

\begin{table}[H]
	\caption{Checklist de configuración TLS del servidor XMPP}
	\label{tab:checklist-tls}
	\centering
	\footnotesize
	\renewcommand{\arraystretch}{1.25}
	\begin{tabularx}{\textwidth}{@{}>{\raggedright\arraybackslash}p{0.34\textwidth} >{\raggedright\arraybackslash}p{0.14\textwidth} Y@{}}
		\toprule
		\textbf{Verificación} & \textbf{Resultado} & \textbf{Evidencia / observación} \\
		\midrule
		Vigencia del certificado & & \\
		Coincidencia del nombre (CN/SAN) & & \\
		Cadena de confianza completa & & \\
		Versión de TLS negociada & & \\
		Suites de cifrado permitidas & & \\
		Interfaz administrativa restringida & & \\
		Cifrado de extremo a extremo (E2EE) & & \\
		Errores de negociación en los registros & & \\
		\bottomrule
	\end{tabularx}
\end{table}

\registrohallazgo{Caso 10}
